% LInA · Artigo acadêmico (entrada, meio, finalização).
% Compile com XeLaTeX. Só comandos de lina.cls: nada de \vspace, \fontsize, \color, \hspace ou \newpage aqui.
\documentclass[artigo]{lina}

\titulocurto{Sobrenome et al. · Título abreviado do artigo}
\cabecalhoprimeira{Cadernos LInA · v. 2 · n. 1 · 2026 · p. 1–18}
\siglarodape{DOI 10.0000/cl.2026.0001}

\begin{document}

% ---------- Entrada ----------
\rotulo{Artigo}
\titulo{Título do artigo em até três linhas, com subtítulo após dois-pontos se necessário}
\autores{\destaque{Nome Sobrenome}¹ · \destaque{Nome Sobrenome}¹ · \destaque{Nome Sobrenome}²}
\afiliacao{¹ Universidade Federal do Ceará · ² Instituição}
\begin{resumo}
  Resumo de até 250 palavras, em parágrafo único: objetivo, método, resultados e~conclusão. O resumo ocupa a largura total e é separado do corpo por dois fios areia.
  \palavraschave{três a cinco termos; separados por ponto e vírgula}
\end{resumo}
\begin{colunas}
\section{1 · Introdução}
O corpo do artigo corre em duas colunas de 79 mm com 8 mm de entrecoluna. Seções numeradas em Bold no mesmo corpo do texto; não há recuo de primeira linha, e os parágrafos se separam por meia linha.

\columnbreak
Citações seguem a norma da publicação. Notas vão ao pé da coluna em 9 pt. Figuras ocupam uma coluna ou a largura total e sempre trazem legenda e fonte.
\end{colunas}

\paginademo

% ---------- Meio ----------
\begin{colunas}
\section{2 · Método}
Descreva o recorte, os dados e os procedimentos em ordem de execução. Quando o método se apoiar em ferramenta própria, nomeie a versão e o repositório em mono no próprio texto.

As subseções usam numeração decimal (2.1, 2.2) em Medium, sem quebra de coluna antes delas.
\end{colunas}

\figuramarcadora[260]
\legenda{Figuras de largura total interrompem as duas colunas e voltam a elas logo abaixo. \fonte{Fonte: LInA, 2026.}}

\begin{colunas}
\section{3 · Resultados}
Apresente os resultados na mesma ordem das perguntas da introdução. Tabelas trazem título acima e fonte abaixo; figuras, legenda abaixo.

\columnbreak
Números em tabela vão em mono, alinhados, com a mesma quantidade de casas decimais por coluna.
\end{colunas}

\titulotabela{Título da tabela acima dela, em uma linha.}
\begin{tabela}{YW{122}W{110}}
  Recorte & Elementos & Relações \\
  Centro & \numero{—} & \numero{—} \\
  Benfica & \numero{—} & \numero{—} \\
  Aldeota & \numero{—} & \numero{—} \\
\end{tabela}

\paginademo

% ---------- Finalização ----------
\begin{colunas}
\section{5 · Conclusão}
Retome o objetivo em uma frase, responda a ele com os resultados e indique o alcance e os limites do trabalho. Termine com os próximos passos da pesquisa.
\end{colunas}

\secaosolta{Agradecimentos}
\textomenor{Agências de fomento, número do processo e pessoas que contribuíram sem autoria.}

\secaosolta{Referências}
\begin{referencias}
BERTALANFFY, Ludwig von. \destaque{Teoria geral dos sistemas}. Petrópolis: Vozes, 1975.

SOBRENOME, Nome. Título do capítulo. \destaque{In}: SOBRENOME, Nome (org.). \destaque{Título do livro}. Cidade: Editora, ano. p. 00–00.

SOBRENOME, Nome; SOBRENOME, Nome. Título do artigo. \destaque{Periódico}, v. 0, n. 0, p. 00–00, ano. DOI: \doi{10.0000/0000}.

\columnbreak
VIEIRA, Jorge de Albuquerque. Organização e sistemas. \destaque{Informática na educação: teoria \& prática}, Porto Alegre, v. 3, n. 1, 2000. DOI: \doi{10.22456/1982-1654.6363}.

SOBRENOME, Nome. \destaque{Título da tese}. Tese (Doutorado em Arquitetura) — Universidade, Cidade, ano.
\end{referencias}

\vfill
\historico
  {\rotulometa{Correspondência}\newline Nome Sobrenome\newline \doi{nome@ufc.br}}
  {\rotulometa{Histórico}\newline \doi{recebido 12/02/2026}\newline \doi{aceito 30/04/2026}}
  {\rotulometa{Licença}\newline CC BY 4.0. Permitida a reprodução com citação da fonte.}

\end{document}
